\chapter{Deep-model scaffolds (not executed on CPU)}
\label{ch:deep-stubs}

\section{Motivation}

Three deep models complete the method coverage and are implemented in
\code{analysis/scripts/deep\_models.py}, guarded behind the
\code{RUN\_DEEP=1} environment flag. None are trained in the current
iteration because no GPU allocation is available; the code is paper-ready
so a single GPU session can produce the missing results.

\section{Stimulus reconstruction CNN}

\paragraph{Input / output.} Input is a 32-ch $\times$ 2-s EEG window
at \SI{64}{\hertz} (shape $(B, 32, 128)$). Output is the reconstructed
28-band gammatone envelope at the same sampling rate (shape
$(B, 28, 128)$).

\paragraph{Architecture.} A three-layer temporal-spatial 1-D CNN:
spatial $1 \times 1$ convolution to 64 channels, two temporal
convolutions (kernel 5) to 128 channels with GELU activation, and a
$1 \times 1$ head projecting to 28 output bands.

\paragraph{Training.} MSE loss between predicted and true attended
envelope on non-overlapping 2-s windows. \code{AdamW}, lr $= 10^{-3}$,
weight decay $10^{-4}$, 50 epochs. Per-subject fine-tuning from a
LOSO pre-trained base.

\paragraph{Evaluation.}
\begin{itemize}
  \item \textbf{Primary}: STOI between reconstructed and attended
        envelope vs.\ reconstructed and unattended envelope; the AAD
        decision is $\text{STOI}(rec, att) > \text{STOI}(rec, una)$.
  \item \textbf{Secondary}: Pearson $r$ with attended vs.\ unattended
        envelope band-by-band.
  \item \textbf{Expected pooled accuracy}: $0.60$--$0.70$ on the
        hemisphere task at $\SI{2}{\second}$ windows, based on the
        KU Leuven CNN-reconstruction baseline (Geirnaert et al.\ 2020
        report $0.65$ at $\SI{2}{\second}$ on a similar stereo paradigm).
\end{itemize}

\section{Foundation-model probe (LaBraM / BIOT)}

\paragraph{Architecture.} Load a pretrained EEG foundation model
(\code{BAAI/LaBraM-Base}, BIOT, or Neuro-GPT variants) with frozen
weights, mean-pool the last hidden state across time, and attach a
single linear classifier for 4-way AAD.

\paragraph{Motivation.} Foundation-model probes are a 2026 publication
staple: they produce a \emph{transfer} baseline that positions the
dataset within the broader EEG pre-training literature without
committing to a particular training recipe. Even a simple linear
probe on a frozen LaBraM gives a reviewer-ready ``does this dataset
benefit from cross-cohort pre-training?'' answer.

\paragraph{Expected pooled accuracy}: hemisphere $\geq 0.70$, 4-class
$\geq 0.35$ — within 5\,pp of the best bespoke model. A probe that
\emph{beats} the bespoke model would be the single strongest result
in the paper.

\section{Cross-modal attention (EEG $\times$ gaze dual-stream)}

\paragraph{Architecture.} A 4-head, 4-layer transformer encoder that
takes EEG tokens (\SI{2}{\second} CNN patches) and gaze tokens
(per-\SI{0.5}{\second} gaze-summary features), linearly embeds each
to a shared 128-D token space, concatenates the two token streams, and
applies joint self-attention before a mean-pool + linear head.

\paragraph{Motivation.} This addresses the fusion ceiling identified in
Chapter~\ref{ch:results}: late / stacked / early fusion all reach
$\sim 0.77$ because the linear combination rule cannot model the
per-trial interaction between gaze fixation and EEG envelope tracking.
Cross-modal attention lets the network learn that interaction.

\paragraph{Expected gain}: $+5$--$10$\,pp over gaze-only on the
hemisphere task. The per-subject interaction is strongest for the
mid-quality tier (5 subjects in Chapter~\ref{ch:indiv-diff}) whose
gaze is only moderately reliable — exactly the subjects for whom
gaze + EEG should compound.

\section{Runtime estimate}

\begin{table}[ht]
\centering
\begin{tabular}{lrr}
\toprule
Model & GPU wall-clock (per subject) & Total (16 subj, single A100) \\
\midrule
StimulusReconstructionCNN   & $\sim 15$ min & $\sim 4$\,h \\
FoundationEEGProbe (frozen) & $\sim 8$ min  & $\sim 2$\,h \\
CrossModalAttention         & $\sim 20$ min & $\sim 5.5$\,h \\
\bottomrule
\end{tabular}
\end{table}

\noindent
Total GPU-hours for the full Stage-C sweep: $\sim 12$\,h on a single
A100, or one overnight run on a PAS2301 GPU node.

\section{Integration with the benchmark}

Each deep model writes its per-subject parquet into
\code{results/deep/\{cnn\_recon,probe,xmod\}/s*.parquet} using the same
schema as the existing classifier scripts
(\code{subject, task, classifier, chance, acc}). Re-running
\code{aggregate\_benchmark\_leaderboard.py} then folds the deep models
into Chapter~\ref{ch:benchmark}'s pooled tables automatically, with
bootstrap CIs computed the same way.
